\newpage
\begin{theorem}[O punkcie stałym rekursji]
    Dla totalnej funkcji rekurencyjnej \(\sigma: \natural \to \natural\) istnieje konstruktywnie \(x_0 \in \natural\) takie, że 
    \[ f_{x_0}(x) = f_{\sigma(x_0)}(x), \]
    czyli obie zatrzymają się i wyniki będą równe.
\end{theorem}
Oznacza to, że dla każdej transformacji \(\sigma\) istnieje maszyna, która nie podlega transformacji i możemy ją skonstruować.
\begin{proof}
    Definiujemy \(G\), konstruktora maszyn:
    \begin{enumerate}
    \singlespacing
        \item Wczytaj \(i\).
        \item Skonstruuj maszynę \(M^{(i)}\) (opis poniżej).
        \item Oblicz indeks \(M^{(i)}\).
    \end{enumerate}
    Maszyna \(M^{(i)}\):
    \begin{enumerate}
    \singlespacing
        \item Wczytaj \(x\).
        \item Skonstruuj \(i\)-tą maszynę \(M_i\).
        \item Oblicz \(k = f_i(i)\), symulując \(M_i\) na wejściu \(i\) (może się zapętlić).
        \item Skonstruuj \(k\)-tą maszynę \(M_k\).
        \item Oblicz i wypisz \(f_k(x)\) (może się zapętlić).
    \end{enumerate}
    
    Jeśli maszyna \(G\) opisuje pewną funkcję totalną \(g: \natural \to \natural\), to
    \[ f_{g(i)}(x) = f_{f_i(i)}(x) \]
    Ustalając \(j\) równego indeksowi funkcji \(\sigma \circ g \) oraz \(x_0 = g(j)\), otrzymujemy:
    \[ f_{x_0}(x) = f_{g(j)}(x) = f_{\sigma(g(j))}(x) = f_{\sigma(x_0)}(x) \]
\end{proof}